%   amsmacros.tex
%
%   Macros (full length headings e.g. for papers/thesis or handout notes)
%   Theorem styles for amsart/amsbook document classes
%
%   Sean McGarraghy
%
%   Change Control:
%   Date      Ver  Reason
%   --------  ---  --------------------------------------------------------------
%   09/01/03  0.1  Begun
%

%   Macros (heading names for papers/thesis or handout notes)
%   Included in other macro files
%
%   Sean McGarraghy
%
%   Change Control:
%   Date      Ver  Reason
%   --------  ---  --------------------------------------------------------------
%   09/01/03  0.1  Begun
%

\newcommand*{\Ax}  {Axiom~}
\newcommand*{\Df}  {Definition~}
\newcommand*{\Th}  {Theorem~}
\newcommand*{\Co}  {Corollary~}
\renewcommand*{\Pr}{Proposition~}
\newcommand*{\Lm}  {Lemma~}
\newcommand*{\Fa}  {Fact~}
\newcommand*{\Qn}  {Question~}
\newcommand*{\Cnj} {Conjecture~}
\newcommand*{\Pb}  {Problem~}
\newcommand*{\Ntn} {Notation~}
\newcommand*{\No}  {Note~}
\newcommand*{\Rk}  {Remark~}
\newcommand*{\Cm}  {Comment~}
%\newcommand*{\Cas} {Case~}
%\newcommand*{\Cs}  {Case~}
\newcommand*{\In}  {Introduction~}
\newcommand*{\Ex}  {Example~}
\newcommand*{\Exe} {Exercise~}
\newcommand*{\Eq}  {Equation~}
\newcommand*{\dfn} {definition}

\newcommand*{\Ch}  {Chapter~}

\newcommand*{\Fi}  {Figure~}
\newcommand*{\Ta}  {Table~}
\newcommand*{\Al}  {Algorithm~}



%\usepackage{amsthm}
\theoremstyle{plain}

\newtheorem{thm}{Theorem}[\numberbychaporsect]
%\newtheorem{lemma}[thm]{Lemma} % commented because beamer has defined it
\newtheorem{lem}[thm]{Lemma}
\newtheorem{propn}[thm]{Proposition}
\newtheorem{pr}[thm]{Proposition}
\newtheorem{prp}[thm]{Propn.}
\newtheorem{cor}[thm]{Corollary}
% \newtheorem{fact}[thm]{Fact} % commented because beamer has defined it
\newtheorem{fct}[thm]{Fact}
\newtheorem{conj}[thm]{Conjecture}
\newtheorem{qn}[thm]{Question}
\newtheorem{probl}[thm]{Problem}
%\newtheorem{claim}{Claim}[]
\newtheorem{prn}[thm]{Principle}

\theoremstyle{definition}

\newtheorem{dmy}[thm]{}  % Dummy theorem name for numbering only
\newtheorem{dm}{}        % Dummy theorem name for numbering only, independent of section
\newtheorem{defn}[thm]{Definition}
%\newtheorem{ex}[thm]{Example}
\newtheorem{exm}[thm]{Example}
\newtheorem{exe}[thm]{Exercise}

\theoremstyle{remark}

\newtheorem{rmk}[thm]{Remark}
\newtheorem{notn}[thm]{Notation}
%\newtheorem{note}[thm]{Note}  % comment if using beamer
\newtheorem{nte}[thm]{Note}
\newtheorem{cmt}[thm]{Comment}
%\newtheorem{case}[thm]{Case}
\newtheorem{intr}[thm]{Introduction}
\newtheorem{intro}[thm]{Introduction}

%\newenvironment{pf}{\vspace{-4pt}\noindent\begin{proof}}{\end{proof}\vspace{2pt}}
%\newenvironment{pf}{\vspace{-4pt}\noindent\begin{proof}}{\end{proof}}
\newenvironment{pf}{\vspace{-1.3em}\noindent\begin{proof}}{\end{proof}}
\newenvironment{eq}{\begin{eqnarray*}}{\end{eqnarray*}}
\newenvironment{ex}{\begin{exm}}{\exEnd\end{exm}}

\usepackage{mathtools}
%\DeclarePairedDelimiter\bra{\langle}{\rvert}
%\DeclarePairedDelimiter\ket{\lvert}{\rangle}
%\DeclarePairedDelimiterX\braket[2]{\langle}{\rangle}{#1 \delimsize\vert #2}

\newcommand*{\hilb}{\mathcal{B}^{\otimes N}}

\usepackage{braket}
\usepackage{hyperref}

\usepackage{graphicx,epsfig,tikz}
\usepackage{hyperref}
\usepackage{caption}
\usepackage{wrapfig}
%\lstset{language=Python}
%\lstset{ 
%  backgroundcolor=\color{white},
%  basicstyle=\footnotesize
%}



%\usepackage{braket}
%\usepackage{hyperref}
\usepackage{bbm}

%\usepackage{braket}
%\usepackage{hyperref}